% Options for packages loaded elsewhere
\PassOptionsToPackage{unicode}{hyperref}
\PassOptionsToPackage{hyphens}{url}
\PassOptionsToPackage{dvipsnames,svgnames,x11names}{xcolor}
%
\documentclass[
  10pt,
]{article}
\usepackage{amsmath,amssymb}
\usepackage{iftex}
\ifPDFTeX
  \usepackage[T1]{fontenc}
  \usepackage[utf8]{inputenc}
  \usepackage{textcomp} % provide euro and other symbols
\else % if luatex or xetex
  \usepackage{unicode-math} % this also loads fontspec
  \defaultfontfeatures{Scale=MatchLowercase}
  \defaultfontfeatures[\rmfamily]{Ligatures=TeX,Scale=1}
\fi
\usepackage{lmodern}
\ifPDFTeX\else
  % xetex/luatex font selection
\fi
% Use upquote if available, for straight quotes in verbatim environments
\IfFileExists{upquote.sty}{\usepackage{upquote}}{}
\IfFileExists{microtype.sty}{% use microtype if available
  \usepackage[]{microtype}
  \UseMicrotypeSet[protrusion]{basicmath} % disable protrusion for tt fonts
}{}
\makeatletter
\@ifundefined{KOMAClassName}{% if non-KOMA class
  \IfFileExists{parskip.sty}{%
    \usepackage{parskip}
  }{% else
    \setlength{\parindent}{0pt}
    \setlength{\parskip}{6pt plus 2pt minus 1pt}}
}{% if KOMA class
  \KOMAoptions{parskip=half}}
\makeatother
\usepackage{xcolor}
\usepackage[margin=1in]{geometry}
\usepackage{longtable,booktabs,array}
\usepackage{calc} % for calculating minipage widths
% Correct order of tables after \paragraph or \subparagraph
\usepackage{etoolbox}
\makeatletter
\patchcmd\longtable{\par}{\if@noskipsec\mbox{}\fi\par}{}{}
\makeatother
% Allow footnotes in longtable head/foot
\IfFileExists{footnotehyper.sty}{\usepackage{footnotehyper}}{\usepackage{footnote}}
\makesavenoteenv{longtable}
\usepackage{graphicx}
\makeatletter
\def\maxwidth{\ifdim\Gin@nat@width>\linewidth\linewidth\else\Gin@nat@width\fi}
\def\maxheight{\ifdim\Gin@nat@height>\textheight\textheight\else\Gin@nat@height\fi}
\makeatother
% Scale images if necessary, so that they will not overflow the page
% margins by default, and it is still possible to overwrite the defaults
% using explicit options in \includegraphics[width, height, ...]{}
\setkeys{Gin}{width=\maxwidth,height=\maxheight,keepaspectratio}
% Set default figure placement to htbp
\makeatletter
\def\fps@figure{htbp}
\makeatother
\setlength{\emergencystretch}{3em} % prevent overfull lines
\providecommand{\tightlist}{%
  \setlength{\itemsep}{0pt}\setlength{\parskip}{0pt}}
\setcounter{secnumdepth}{5}
\DeclareUnicodeCharacter{2212}{\ensuremath{-}}
\DeclareUnicodeCharacter{2248}{\ensuremath{\approx}}
\DeclareUnicodeCharacter{2264}{\ensuremath{\le}}
\DeclareUnicodeCharacter{2265}{\ensuremath{\ge}}
\DeclareUnicodeCharacter{00D7}{\ensuremath{\times}}
\ifLuaTeX
  \usepackage{selnolig}  % disable illegal ligatures
\fi
\usepackage[round]{natbib}
\bibliographystyle{plainnat}
\IfFileExists{bookmark.sty}{\usepackage{bookmark}}{\usepackage{hyperref}}
\IfFileExists{xurl.sty}{\usepackage{xurl}}{} % add URL line breaks if available
\urlstyle{same}
\hypersetup{
  pdftitle={Lineage Is Memory: Build on the Best Parent and Let It Accumulate What Feedback Reveals},
  pdfauthor={Anonymous},
  colorlinks=true,
  linkcolor={Maroon},
  filecolor={Maroon},
  citecolor={Blue},
  urlcolor={Blue},
  pdfcreator={LaTeX via pandoc}}

\title{Lineage Is Memory: Build on the Best Parent and Let It Accumulate
What Feedback Reveals}
\author{Anonymous}
\date{}

\begin{document}
\maketitle
\begin{abstract}
LLM-guided optimizers such as GEPA improve a prompt by showing a
proposer model a few failing examples and asking it to rewrite. Much
recent work tries to make that step better. It shows the proposer more
of the error distribution through full-set clustering, targeted
minibatches or agents that read every trace, or it pays for a stronger
proposer. We show that a simpler choice matters more: whether each
rewrite builds on the current best candidate (lineage) or starts again
from the seed.

In a controlled grid at equal task-call budget, building on the best
parent adds +0.238 held-out accuracy on seven task families with
gpt-4.1-mini and +0.222 with claude-haiku-5-5. This is larger than any
change to what the proposer is shown. Switching the proposer from
gpt-4.1-mini to gpt-4.1 adds nothing (−0.019 {[}−0.053, +0.014{]}),
while lineage still adds +0.380 under the stronger proposer. Lineage
also absorbs most of what wider reflection buys. A full-failure summary
that keeps the expected answer is the best display when rewriting the
seed (+0.156 over the first four failures) and ties raw failures under
lineage (−0.018 {[}−0.053, +0.010{]}). The same pattern holds on Haiku.

The mechanism is memory. A child inherits the rules its parent already
learned, so facts revealed by feedback accumulate across rounds instead
of having to fit in one prompt. On the Banking77 intent benchmark, the
reference loop at first lost to GEPA by 0.29. Logs traced the gap to one
sentence in GEPA's reflection prompt that asks the proposer to copy task
facts from the feedback into the instruction. Adding a sentence of that
kind, a fix designed after the diagnosis and pre-registered before
running, raises our loop from 0.372 to 0.856. That is above GEPA's 0.664
by +0.192 {[}+0.097, +0.289{]}, at 2.3× fewer tokens. With lineage, the
best prompts name 11--12 of the 12 intents; rewriting the seed with the
same sentence they name 4, and score 0.317. On BIG-Bench Hard error
detection, where there is no vocabulary to carry, the fix does not help
(−0.061), and the methods tie.

Every pre-registered test that failed is reported in §4.8. We release
the harness.
\end{abstract}

\hypertarget{introduction}{%
\section{Introduction}\label{introduction}}

Reflective prompt optimizers have become the default way to tune LLM
programs. GEPA \citep{a2025gepa} runs the current program on a small
minibatch, shows a reflection model the inputs, outputs and feedback,
and accepts the rewrite if it scores higher; Pareto-front parent
selection keeps the search building on what already works. The same
loop, with scores instead of text feedback, drives OPRO
\citep{yang2023large}, EvoPrompt \citep{guo2023evoprompt} and the
program-evolution systems behind FunSearch
\citep{romeraparedes2023mathematical} and AlphaEvolve
\citep{alexander2025alphaevolve}.

Most attempts to improve these loops target the reflection step. ESPO
\citep{liu2026espo} argues from a coupon-collector bound that a
minibatch of three failures covers a small fraction of the error modes
per round, and replaces it with full-set error clustering. RPT
\citep{bayat2026reflective} gives the optimizer a tool that clusters
critiques over the whole training set. AdaGEPA \citep{chen2026adagepa}
swaps minibatch slots toward identified weaknesses. CASD
\citep{singh2026coding} lets a coding agent read the full trace corpus.
The common alternative is to pay for a stronger proposer. All of these
were evaluated inside lineage-based search, where each candidate
rewrites the current best. None turns lineage off, so none can say how
much of the loop's progress comes from what the proposer sees and how
much from what it inherits.

We treat an LLM-guided optimizer as a loop with three knobs. \textbf{P}
(lineage) is which candidate the proposer rewrites. \textbf{F}
(disclosure) is what it is shown about failures. \textbf{A} (scheduling)
is how candidates are evaluated and stopped. We cross P and F in a
single controlled grid at equal task-call budget. We measure coverage
per failure category rather than only aggregate score, and we decompose
cost by call type. The tasks are seven families: five constructed so
that the needed rules can only be learned from feedback, a free-text
task scored by an LLM judge, and ScoNe \citep{she2023scone}. We add two
public benchmarks, Banking77 \citep{casanueva2020efficient} and
BIG-Bench Hard salient-translation-error detection
\citep{suzgun2022challenging}. We replicate on a second model family and
test a stronger proposer.

Our findings:

\begin{enumerate}
\def\labelenumi{\arabic{enumi}.}
\tightlist
\item
  \textbf{Lineage is the main lever, and it is free.} Building on the
  best parent adds +0.238 {[}+0.224, +0.252{]} with gpt-4.1-mini and
  +0.222 {[}+0.204, +0.239{]} with claude-haiku-5-5. On the public
  benchmarks it adds +0.088 {[}+0.045, +0.130{]}. With gpt-4.1 as
  proposer, it adds +0.380 {[}+0.337, +0.425{]}, while the stronger
  proposer itself adds nothing over gpt-4.1-mini (§4.1--4.2).
\item
  \textbf{Lineage is memory, and it can be fed.} On Banking77, the facts
  the task model needs (the intent names) are visible only in feedback.
  With lineage, an instruction to carry those facts into the prompt
  makes them accumulate: the best prompts name 11--12 of 12 intents,
  against 4 when every rewrite starts from the seed. Held-out accuracy
  rises from 0.372 to 0.856, above GEPA's 0.664 at 2.3× fewer tokens
  (§4.3). The fix was designed after a diagnosis of where GEPA won; it
  was pre-registered before it ran and is tested on two tasks.
\item
  \textbf{Lineage makes better failure samples largely redundant.} A
  diverse sample and a full-failure summary that keeps the expected
  answer both help when rewriting the seed and add nothing under
  lineage. The interactions are +0.068 and +0.174 on gpt-4.1-mini, and
  +0.122 and +0.187 on Haiku. GEPA-style random minibatches are worse
  than raw first-k failures under lineage (−0.061 {[}−0.086, −0.035{]}).
  Without lineage, partial disclosure can even push coverage of unshown
  failures below showing nothing (§4.4--4.6).
\item
  \textbf{Most of GEPA's extra cost is validation, not reflection.} At
  matched task calls, building on the best parent with early stopping
  needs 2.6--3.3× fewer tokens to solve than GEPA where both solve.
  Reflection is 2--15\% of GEPA's tokens; on the public benchmarks it is
  8--12\% (§4.7).
\end{enumerate}

The practical upshot is a default and a reporting rule. The default:
build on the best parent, show raw failures with their expected answers,
and tell the proposer to keep what the feedback reveals. The rule: a
method that changes what the optimizer sees should be evaluated with
lineage switched off as well as on, or its gains cannot be attributed.

\hypertarget{related-work}{%
\section{Related work}\label{related-work}}

\textbf{LLM-guided program evolution.} FunSearch
\citep{romeraparedes2023mathematical} and AlphaEvolve
\citep{alexander2025alphaevolve} have a proposer model rewrite programs
drawn from an archive of high scorers. Islands, archives or Pareto
fronts decide which parent is rewritten next; ImprovEvolve
\citep{kravatskiy2026improvevolve} continues this line. In this family
the gain is attributed to evolution itself, meaning lineage plus
selection, and what the evaluator tells the proposer is treated as an
engineering detail. The closest methodological study, Evolution or
Illusion? \citep{oved2026evolution}, sweeps how a fixed proposal budget
is split between seeds and iterations and finds that strategy rankings
flip with budget. It does not vary what the proposer sees, does not
count tokens, and treats best-of-N only as a reference line. We hold the
budget fixed and vary the two factors it leaves alone: lineage and
disclosure.

\textbf{Reflective prompt optimization.} ProTeGi
\citep{pryzant2023automatic} and TextGrad
\citep{yuksekgonul2024textgrad} turn failing examples into textual
``gradients''; OPRO \citep{yang2023large} shows the optimizer only a
trajectory of scores. GEPA \citep{a2025gepa} combines reflection on a
random minibatch of three examples with Pareto parent selection, and it
has become the default baseline. Its reflection prompt asks the model to
``identify all niche and domain specific factual information about the
task and include it in the instruction''; §4.3 shows that this sentence,
not the minibatch or the Pareto front, accounts for GEPA's advantage
over our loop on Banking77. Later analyses question what drives GEPA.
\citep{lumer2026reflective} find that the \emph{form} of feedback
matters: per-dimension feedback transfers, a scalar does not. They
attribute validation-to-held-out gaps to a winner's curse under noisy
judges. \citep{giummarra2026understanding} report that GEPA helps on
MATH-500 but not on HotpotQA with local models, and that it raises token
use. \citep{ma2024large} argue that LLM optimizers rely mostly on their
own prior rather than on the error feedback they receive, which matches
our null on ScoNe and value labels (§4.5).

\textbf{Wider reflection scope (2026).} A cluster of 2026 papers argues
that a few random failures under-cover the error distribution. ESPO
\citep{liu2026espo} states the problem as incomplete error observation.
It bounds coverage with a coupon-collector argument (about 15 rounds for
a 95\% chance of seeing every mode at minibatch three) and replaces the
minibatch with full-set error clustering into 3--7 patterns, followed by
Diversify and Stabilize steps. RPT \citep{bayat2026reflective} gives the
optimizer a diagnosis tool that clusters critiques over the whole
training set. AdaGEPA \citep{chen2026adagepa} swaps at most one
minibatch example per round to target an identified weakness. CASD
\citep{singh2026coding} lets a coding agent read the full corpus and
reports wins over GEPA on three of four benchmarks at lower cost
{[}abstract-level{]}. Contrastive Reflection \citep{koh2026contrastive}
selects failures by decision-tree slices, a +1.4-point gain that lies
inside its bootstrap interval. Each of these is evaluated inside
lineage-based search, usually GEPA itself. None measures per-category
coverage against a no-feedback prior, and none turns lineage off. We
find that the problem they address is real when the proposer rewrites a
fixed seed, and that lineage removes most of it (§4.4, §4.6).

\textbf{Search efficiency.} Explore More, Validate Less
\citep{zhou2026explore} observes that GEPA spends evaluations on
all-correct minibatches that yield no edit, and proposes fewer full
validations. It does not break GEPA's budget down by call type; §4.7
does, and shows that reflection is a small share.

\textbf{Representative examples.} Choosing a small set of examples that
represents a larger one is an old problem in program synthesis
\citep{pu2017selecting} and in-context learning. Our \emph{diverse} arm
borrows that idea; what is new here is measuring its effect inside an
evolution loop and against lineage.

\hypertarget{setup-the-pfa-loop}{%
\section{Setup: the P/F/A loop}\label{setup-the-pfa-loop}}

\textbf{Loop.} Each run starts from a fixed seed: a one-line instruction
for prompt tasks, or a deliberately naive function for code tasks. It
spends a fixed budget of candidate evaluations, 8 for prompts and 12 for
code. In each step a proposer model receives the objective, one
candidate and the evaluator's feedback on it, and returns a rewrite,
which is scored on the search set. The best search-set candidate of the
run is then scored once on a disjoint held-out set, which is never used
for selection. A run stops early when a candidate scores 1.0 on the
search set.

\textbf{P (lineage).} \emph{seed}: every proposal rewrites the seed, so
proposals are independent samples conditioned on the same feedback
(best-of-N). \emph{parent}: every proposal rewrites the current best
candidate, so fixes accumulate.

\textbf{F (disclosure).} The evaluator shows up to four failing cases
with the expected answer. The arms are:

\begin{itemize}
\tightlist
\item
  \emph{none}: score only.
\item
  \emph{first}: the first four failures in dataset order.
\item
  \emph{diverse}: four failures picked by farthest-point selection under
  edit distance on a coarse shape of the input.
\item
  \emph{minibatch}: four examples drawn at random, right or wrong, as in
  GEPA's reflective minibatch.
\item
  \emph{cluster}: all failures are passed to one extra model call that
  groups them into three to seven patterns with one verbatim example
  each, and only that summary is shown. This is a stand-in for the
  full-set diagnosis of ESPO \citep{liu2026espo}, and its tokens are
  charged to the run.
\item
  \emph{clustergold}: as \emph{cluster}, but each pattern's example
  keeps its input and expected answer (§4.6).
\end{itemize}

\textbf{A (scheduling).} Our loop evaluates every candidate on the full
search set and stops at the first perfect score. GEPA 0.0.7
\citep{a2025gepa} is run as published, with minibatch 3, Pareto parent
selection and the same task-model call budget (8 × search-set size). It
is run again with an early-stop hook that ends the run at its first
perfect full-set evaluation.

\textbf{Tasks.} There are seven families in the main grid:

\begin{itemize}
\tightlist
\item
  Five constructed families in which the needed rules cannot be read off
  the task statement: duration and size parsing (code, Python functions
  run in a sandbox), value labels (a closed label set with idiosyncratic
  boundaries), date normalisation, and ticket routing.
\item
  One free-text family, support replies, scored by a gpt-4.1-mini judge
  at temperature 0 against per-example rubrics.
\item
  ScoNe \citep{she2023scone}, using its six negation conditions as
  categories (6 search and 12 held-out pairs per condition,
  label-balanced).
\end{itemize}

Two further public benchmarks were added in a pre-registered extension:

\begin{itemize}
\tightlist
\item
  \textbf{Banking77} \citep{casanueva2020efficient}: 12 of its 77
  intents, 24 search and 36 held-out queries. The seed names no intent
  labels, so the label vocabulary is visible only in feedback.
\item
  \textbf{BIG-Bench Hard salient-translation-error detection}
  \citep{suzgun2022challenging}: 18 search and 36 held-out items. The
  shared instructions and options are moved into the seed, and inputs
  longer than the evaluator's 300-character feedback window are dropped.
\end{itemize}

Every example carries a category label, so we measure coverage (the
share of candidates that get every example of a category right) in
addition to the aggregate score.

\textbf{Models and runs.} In the main grid, gpt-4.1-mini is the task
model, the proposer and the clusterer. Each of the 9 cells (seed-none
and \{seed, parent\} × \{first, diverse, minibatch, cluster\}) is run
with 10 seeds per task. The replication uses claude-haiku-5-5 in all
three roles (§4.1, §4.4). The strong-proposer arm uses gpt-4.1 as
proposer and clusterer with gpt-4.1-mini as task model (§4.2). Intervals
are 95\% bootstrap intervals over runs; contrasts that average over
tasks resample runs within each task. Each protocol and its success
threshold were registered before its runs (Appendix A).

\hypertarget{results}{%
\section{Results}\label{results}}

All numbers are held-out best scores: the mean over 10 seeds with a 95\%
bootstrap interval. Cross-task contrasts weight tasks equally
(\texttt{results/final-7tasks.txt}, \texttt{results/round4.txt},
\texttt{results/round5.txt}). Post-hoc analyses are labelled as such.

\begin{figure}
\centering
\includegraphics{../figs/fig1-lineage-disclosure.png}
\caption{Held-out best by disclosure, without (seed) and with (parent)
lineage.}
\end{figure}

\textbf{Figure 1.} Held-out best by disclosure scheme on the seven main
tasks, seed row (orange) versus parent row (blue); bands are 95\%
intervals over 10 seeds. ``none'' is shared by both rows.

\hypertarget{lineage-is-the-main-lever}{%
\subsection{Lineage is the main lever}\label{lineage-is-the-main-lever}}

Averaged over the four disclosure schemes and seven tasks, building on
the best parent raises the held-out score by +0.238 {[}+0.224,
+0.252{]}. Among the four pre-registered schemes, no contrast within a
row exceeds 0.10. The gold-keeping summary of §4.6 widens the seed-row
spread to 0.156, still well below the lineage effect. On the three tasks
with room to move and rules that must be learned from feedback (labels,
routing, replies), every pre-registered parent cell beats every
pre-registered seed cell (Fig. 1). The one crossing involves the
sensitivity arm, where seed-row clustergold on replies (0.794) exceeds
parent-row cluster (0.519). Dates and ScoNe are flat: dates is solved
with no feedback at all (seed-none 1.000), and ScoNe sits at 0.89--0.92
in every cell.

\textbf{Second model family.} The effect carries over to
claude-haiku-5-5 as task model, proposer and clusterer. On five tasks
(labels, routing, replies, duration, size), averaged over three
disclosure columns, lineage adds +0.222 {[}+0.204, +0.239{]}
(pre-registered H1).

\textbf{Public benchmarks.} On Banking77 and BIG-Bench Hard together,
averaged over four disclosure columns, lineage adds +0.088 {[}+0.045,
+0.130{]} (pre-registered P1). The effect comes from Banking77 (+0.158
{[}+0.077, +0.241{]}). On BBH it is +0.017 {[}−0.006, +0.039{]}: every
cell there lies between 0.61 and 0.68, and seed-none already reaches
0.556.

\hypertarget{a-stronger-proposer-does-not-substitute-for-lineage}{%
\subsection{A stronger proposer does not substitute for
lineage}\label{a-stronger-proposer-does-not-substitute-for-lineage}}

We replaced the proposer and clusterer with gpt-4.1, kept gpt-4.1-mini
as task model and judge, and ran five cells (seed-none and \{seed,
parent\} × \{first, clustergold\}) on labels, routing and replies. All
three pre-registered tests hold:

\begin{itemize}
\tightlist
\item
  Lineage adds +0.380 {[}+0.337, +0.425{]} (S1); with the mini proposer
  the same gap is +0.365 {[}+0.316, +0.413{]}.
\item
  The gold-keeping summary helps far more without lineage:
  I\_clustergold = +0.294 {[}+0.201, +0.381{]} (S2).
\item
  Building on the best parent, it is slightly worse than first-k: −0.057
  {[}−0.101, −0.016{]} (S3).
\end{itemize}

Cell by cell, the stronger proposer does not beat the mini one
(secondary, not pre-registered as a test):

\begin{longtable}[]{@{}ll@{}}
\toprule\noalign{}
Cell & gpt-4.1 − gpt-4.1-mini proposer \\
\midrule\noalign{}
\endhead
\bottomrule\noalign{}
\endlastfoot
parent-first & −0.019 {[}−0.053, +0.014{]} \\
parent-clustergold & −0.025 {[}−0.087, +0.057{]} \\
seed-clustergold & −0.069 {[}−0.170, +0.033{]} \\
seed-none & +0.002 \\
\end{longtable}

At these budgets, buying a stronger proposer bought nothing, while
choosing to build on the best parent was worth about 0.38.

\hypertarget{lineage-is-memory-carrying-what-feedback-reveals}{%
\subsection{Lineage is memory: carrying what feedback
reveals}\label{lineage-is-memory-carrying-what-feedback-reveals}}

\textbf{The gap.} On Banking77, GEPA with minibatch 3 reached 0.664
held-out at the same task-call budget, against 0.372 for our loop with
lineage and first-k failures. The difference is +0.292 {[}+0.111,
+0.464{]} (+0.289 with early stopping). On BBH the two tied (−0.019
{[}−0.050, +0.014{]}). No run of either method solved the search set on
either task, so the pre-registered token-ratio test (R4) could not be
evaluated there.

\textbf{Diagnosis (post hoc, from the logs).} Two observations explain
the gap.

\begin{itemize}
\tightlist
\item
  \emph{Stalling at zero.} The Banking77 seed names no labels and scores
  0. First-generation rewrites mostly fixed the output format
  (snake\_case) and still scored 0, so the best parent stayed the seed.
  In two of ten runs all eight candidates scored 0, and the final prompt
  was the 36-character seed. Lineage cannot help on a flat plateau.
\item
  \emph{Fact transfer.} The feedback revealed 12 intent names. Our best
  prompts named {[}6, 11, 4, 11, 7, 9, 4, 4, 0, 0{]} of them, which
  tracks their held-out scores {[}0.47, 0.69, 0.31, 0.75, 0.39, 0.58,
  0.22, 0.31, 0, 0{]}. GEPA's best prompts named a median of 9. GEPA
  0.0.7's reflection prompt explicitly asks the model to copy ``niche
  and domain specific factual information'' from the feedback into the
  instruction; our proposer's objective had no such sentence. BBH has no
  vocabulary to carry, and there the methods tie.
\end{itemize}

\textbf{Fix (pre-registered before running).} We appended one sentence
to the proposer's objective and changed nothing else: \emph{``The task
model sees only this prompt, never the feedback: copy every
task-specific fact the feedback reveals (for example the exact expected
labels or formats) into the prompt.''} We ran seed-first and
parent-first with it on both public tasks, 10 seeds each.

\begin{longtable}[]{@{}
  >{\raggedright\arraybackslash}p{(\columnwidth - 8\tabcolsep) * \real{0.2000}}
  >{\raggedright\arraybackslash}p{(\columnwidth - 8\tabcolsep) * \real{0.2000}}
  >{\raggedright\arraybackslash}p{(\columnwidth - 8\tabcolsep) * \real{0.2000}}
  >{\raggedright\arraybackslash}p{(\columnwidth - 8\tabcolsep) * \real{0.2000}}
  >{\raggedright\arraybackslash}p{(\columnwidth - 8\tabcolsep) * \real{0.2000}}@{}}
\toprule\noalign{}
\begin{minipage}[b]{\linewidth}\raggedright
task
\end{minipage} & \begin{minipage}[b]{\linewidth}\raggedright
parent-first
\end{minipage} & \begin{minipage}[b]{\linewidth}\raggedright
seed-first + carry
\end{minipage} & \begin{minipage}[b]{\linewidth}\raggedright
parent-first + carry
\end{minipage} & \begin{minipage}[b]{\linewidth}\raggedright
GEPA mb3
\end{minipage} \\
\midrule\noalign{}
\endhead
\bottomrule\noalign{}
\endlastfoot
Banking77 & 0.372 & 0.317 & \textbf{0.856} & 0.664 \\
BBH ste & 0.653 & 0.614 & 0.592 & 0.633 \\
\end{longtable}

On Banking77 the fix adds +0.483 {[}+0.328, +0.644{]} (F1). It overtakes
GEPA by 0.192: GEPA − ours = −0.192 {[}−0.289, −0.097{]} (F2), at 45.8k
against 106.9k tokens per run (2.3×). No run stalls at zero (0/10,
against 2/10 before).

The fix works only with lineage. Rewriting the seed with the same
sentence scores 0.317, and lineage adds +0.539 {[}+0.494, +0.583{]} on
top of it (pooled over both tasks, F4: +0.258 {[}+0.221, +0.296{]}). The
prompts show why. With lineage, the best prompts name 11--12 of the 12
intents in every run. Rewriting the seed, they name exactly 4 in every
run, the intents in the four failures every seed-row round sees. A child
keeps the labels its parent already wrote down and adds those its own
feedback reveals, so the vocabulary accumulates across rounds. A
proposer that restarts from the seed can only write down what one round
shows it.

On BBH, the fix does not help. It lowers parent-first by −0.061
{[}−0.119, −0.006{]}, so the pre-registered no-harm test F3 fails. GEPA
and our fixed loop are within the noise there (+0.042 {[}−0.008,
+0.094{]}). Carrying facts pays when the feedback contains a vocabulary
the task model cannot guess, and it is not a general improvement.

\begin{figure}
\centering
\includegraphics{../figs/fig2-banking77-carry.png}
\caption{Banking77: carrying facts across lineage.}
\end{figure}

\textbf{Figure 2.} Banking77 held-out best (left) and number of the 12
intents named in the best prompt (right). Orange: rewrite the seed;
blue: build on the best parent; grey: GEPA (minibatch 3, same task-call
budget); ``+carry'' adds the one-sentence fact-carrying instruction.
Bars are 95\% intervals over 10 seeds (\texttt{results/fig2.py}).

\hypertarget{lineage-absorbs-most-of-what-a-better-sample-buys-c2}{%
\subsection{Lineage absorbs most of what a better sample buys
(C2)}\label{lineage-absorbs-most-of-what-a-better-sample-buys-c2}}

Without lineage, a diverse sample beats the first four failures by
+0.061 {[}+0.046, +0.077{]}; with lineage the same contrast is −0.007
{[}−0.025, +0.012{]}. The pre-registered interaction is I\_diverse =
+0.068 {[}+0.043, +0.092{]}, so C2 is supported and refutation R1 does
not hold. On claude-haiku-5-5 the interaction is +0.122 {[}+0.088,
+0.158{]} (pre-registered Hd): diverse helps rewriting the seed (+0.072
{[}+0.055, +0.090{]}) and does not help building on the parent (−0.050
{[}−0.082, −0.020{]}). The Haiku first-k cells come from the earlier
replication batch, not interleaved with the diverse cells. Routing is
the exception on both models: on Haiku its interaction is −0.092
{[}−0.142, −0.042{]}, because there diverse is worse than first-k
without lineage.

The interaction for the clustered summary is larger, +0.120 {[}+0.073,
+0.167{]}, but for a different reason: in the seed row the summary helps
a little (+0.038 over first), while in the parent row it is worse
(§4.6). With a summary that keeps the expected answer the interaction
grows to +0.174 {[}+0.127, +0.222{]} (§4.6).

\hypertarget{without-lineage-partial-disclosure-can-crowd-out-unshown-failures-c1}{%
\subsection{Without lineage, partial disclosure can crowd out unshown
failures
(C1)}\label{without-lineage-partial-disclosure-can-crowd-out-unshown-failures-c1}}

In the seed row every proposal sees the seed's failures, so the
categories that are and are not shown are fixed. On the two code tasks,
showing the first four failures lowers the share of proposals that get
an unshown category fully right \emph{below} showing none: duration
−0.055 {[}−0.082, −0.033{]}, size −0.113 {[}−0.144, −0.082{]}. The
intervention on duration confirms the direction: showing only ISO-style
failures lowers unshown coverage by −0.048 {[}−0.079, −0.021{]}, while a
random sample of four raises it by +0.061 {[}−0.003, +0.121{]}. Value
labels show no crowding-out (+0.018 {[}−0.031, +0.067{]}), so C1 holds
on 2 of 3 constructed tasks, meeting the registered threshold. On ScoNe
every seed-row disclosure arm is within ±0.01 of seed-none (R3 holds).
Because ScoNe is near ceiling, this is a boundary condition, not
evidence of no effect.

\hypertarget{wider-reflection-scope-helps-only-without-lineage}{%
\subsection{Wider reflection scope helps only without
lineage}\label{wider-reflection-scope-helps-only-without-lineage}}

In the parent row the ordering is first ≈ diverse \textgreater{}
minibatch ≈ cluster. GEPA-style random minibatches lose −0.061
{[}−0.086, −0.035{]} against first, and the clustered summary loses
−0.082 {[}−0.119, −0.045{]}. The pre-registered prediction that the
clustered summary beats the best alternative in both rows (R2) is
refuted in the opposite direction: −0.044 {[}−0.078, −0.018{]} in the
seed row and −0.095 {[}−0.132, −0.063{]} in the parent row. The replies
deficit has a visible cause. The clustering prompt asks for a
representative wrong answer, and on replies, whose expected outputs are
long rubrics, only 31 of 156 summaries kept the expected answer, against
95--100\% on every other task.

\textbf{A summary that keeps the target (post hoc, registered before
running).} We changed one phrase of the clustering prompt so that each
pattern carries a verbatim failure \emph{with its expected answer}. We
then reran both rows on all seven tasks (10 seeds, same budget). Every
summary kept the expected answer (replies 65 of 65). On replies the
parent row recovers from 0.519 to 0.869 (+0.350 {[}+0.062, +0.612{]}),
so the earlier deficit was information loss, not summarization per se.

Across tasks the picture splits by lineage. Rewriting the seed, the
gold-keeping summary is the best arm we ran: +0.156 {[}+0.122, +0.188{]}
over first-k and +0.095 {[}+0.058, +0.127{]} over diverse. Its gains are
on labels (0.408→0.683), replies (0.294→0.794) and routing
(0.500→0.625). Building on the best parent, it is indistinguishable from
first-k (−0.018 {[}−0.053, +0.010{]}); the interval excludes an average
gain above +0.01, and the contrast is −0.002 without replies. It is also
indistinguishable from diverse (−0.011), still beats random minibatches
(+0.043 {[}+0.005, +0.075{]}), and loses on routing (−0.083 {[}−0.158,
−0.008{]}). Lineage still adds +0.139 {[}+0.096, +0.182{]} on top of the
best summary.

The summary is not free: it costs an extra model call per round, and
median search tokens per run reach up to 2.2× those of first-k (routing,
seed row). Full-set diagnosis as ESPO specifies it is therefore a real
fix for the seed-rewriting regime, and redundant once the search builds
on its best parent. The claude-haiku-5-5 replication reproduces this
split (pre-registered): I\_clustergold = +0.187 {[}+0.152, +0.221{]}
(H2), and under lineage the summary ties first-k at −0.008 {[}−0.017,
+0.002{]} (H3). Under lineage, random minibatches again lose to first-k
(−0.093 {[}−0.123, −0.067{]}).

\hypertarget{where-gepas-tokens-go-c3}{%
\subsection{Where GEPA's tokens go
(C3)}\label{where-gepas-tokens-go-c3}}

At the same task-call budget, evolve with lineage and early stop solves
the search set at least as often as GEPA with early stop on four of five
tasks: labels 7 vs 6, dates 10 vs 10, routing 9 vs 2, replies 10 vs 2.
It solves less often on ScoNe (2 vs 4). Where both solve often enough to
compare, GEPA needs 3.34× {[}2.23, 4.49{]} (labels) and 2.59× {[}2.30,
4.84{]} (dates) the tokens to solve.

Refutation R4 (the gap is explained by GEPA's lack of early stopping) is
rejected. Early stopping changes GEPA's median tokens-to-solve by at
most 6\% on labels, dates and replies; routing has too few solved runs
to say. Only on ScoNe does it halve them (97.8k → 50.6k).

Of GEPA's tokens, reflection is 2--15\%; full re-evaluations take
21--67\% and minibatch checks 27--77\%. On the public benchmarks,
reflection is 8\% (Banking77) and 12\% (BBH), and full re-evaluations
are 74\% and 58\%. GEPA used 2.3--2.8× the tokens per run of our loop
there.

\hypertarget{boundary-conditions-and-robustness}{%
\subsection{Boundary conditions and
robustness}\label{boundary-conditions-and-robustness}}

These are the results that limit the claims above. All pre-registered
tests that failed are listed here.

\begin{itemize}
\tightlist
\item
  \textbf{Disclosure contrasts on the public benchmarks (pre-registered
  P2--P5) do not hold.} The interactions I\_clustergold (+0.017
  {[}−0.119, +0.151{]}) and I\_diverse (−0.014 {[}−0.118, +0.093{]}) are
  not distinguishable from 0. The parent-row tie test P3 (+0.021
  {[}−0.101, +0.143{]}) fails on width: its upper bound exceeds +0.05.
  The minibatch penalty P5 (−0.011 {[}−0.113, +0.090{]}) is not
  distinguishable from 0. On BBH, every disclosure contrast is within
  ±0.04. On Banking77, intervals are about ±0.24 because individual runs
  range from 0 to 0.92. On these two tasks, what the proposer was shown
  mattered less than lineage did, and no disclosure scheme was reliably
  better than first-k.
\item
  \textbf{Carrying facts can cost a little where there are none to
  carry.} On BBH the fix lowers parent-first by −0.061 (F3 fails), and
  under the fix lineage does not help there either (−0.022 {[}−0.086,
  +0.039{]}).
\item
  \textbf{Routing} is the one task where diverse samples hurt without
  lineage (both models), and where clustergold loses under lineage.
\item
  \textbf{ScoNe} is near ceiling (0.89--0.92 in every arm), and GEPA
  solves it more often than our loop (4 vs 2).
\item
  \textbf{Leave-one-task-out (post hoc).} Dropping one main-grid task at
  a time, I\_diverse stays positive (+0.046 to +0.094), and parent-row
  minibatch − first stays negative (−0.042 to −0.069), in all seven
  subsets. The ungrounded clustered deficit does not: without replies,
  parent-row cluster − first is −0.019 {[}−0.045, +0.008{]}. On Haiku,
  H1 and H2 stay positive with any one task dropped, and the H3 upper
  bound stays at or below +0.009.
\item
  \textbf{Noise.} Same-seed reruns of five replies-cluster runs differed
  by up to 0.25, making that cell the noisiest in the grid.
\end{itemize}

\hypertarget{discussion}{%
\section{Discussion}\label{discussion}}

\textbf{Why lineage is memory.} A proposer that rewrites the seed starts
every round from the same failures, so whatever it is shown is the whole
of its information about the task. If the sample is skewed, every round
inherits the skew, and crowding-out (§4.5) is the visible symptom. A
proposer that rewrites the best parent inherits every rule earlier
rounds already learned, and the failures it sees next are the residual
ones, which shrink and shift as search proceeds. Coverage then
accumulates across rounds instead of having to be delivered in one
prompt. This is the coupon-collector argument of \citep{liu2026espo} run
in the other direction: if each round fixes what it saw and keeps it,
the rounds collect the coupons, and the per-round sample does not need
to cover the distribution.

The logs agree (post hoc, first-k cells,
\texttt{results/mechanism\_posthoc.txt}). Over a run, the failures shown
to a seed-row proposer are always the same four, 17--25\% of the search
set. The failures shown to a parent-row proposer cover 54\% (labels),
61\% (routing) and 86\% (replies) of it, and a child keeps 70--91\% of
the examples its parent already got right. Banking77 makes the memory
literal (§4.3). With lineage and an instruction to keep what feedback
reveals, the prompt becomes a growing list of the intents seen so far,
11--12 of 12 by the end of the budget. Without lineage, the same
instruction yields the same 4 intents in every run.

\textbf{Why this matters more than a stronger proposer.} Disclosure and
proposer strength both act on a single round. Lineage acts across
rounds. At a budget of eight evaluations, a proposer that gets one round
slightly better gains less than a loop that keeps every round's gains.
That is what we see: gpt-4.1 in place of gpt-4.1-mini changes the best
score by −0.019, and lineage changes it by +0.38.

\textbf{What this means for practitioners.} If your optimizer builds on
its best candidate, as GEPA, FunSearch-style archives and our loop do,
the cheap defaults are these:

\begin{itemize}
\tightlist
\item
  Show raw failures with their expected answers.
\item
  Tell the proposer that the task model never sees the feedback, so
  anything the feedback reveals must be written into the prompt.
\item
  Spend the budget on more rounds, not on wider reflection or a bigger
  proposer.
\item
  Do not draw a random minibatch of right and wrong examples: under
  lineage, it was the one display choice that hurt in every
  leave-one-task-out subset on gpt-4.1-mini and on Haiku.
\end{itemize}

Full-set diagnosis pays off when the system cannot build on a parent:
parallel best-of-N rewrites of one prompt, one-shot repair from a log of
failures, or a first round before any lineage exists. Whatever the
format, the expected answer must survive it. A summary that kept what
went wrong but not what was wanted lost 0.46 on free-text replies.

\textbf{What this means for evaluating reflection-scope methods.} A
reflection-scope fix measured inside GEPA measures its residual value on
top of lineage. In our loop that residual is near zero against raw
first-k failures but positive against a random minibatch (+0.043
{[}+0.005, +0.075{]}). Gains reported over GEPA's default are therefore
consistent with our results, and may come largely from replacing the
random minibatch, which showing raw failures also does. We did not run
ESPO's full pipeline inside GEPA; this is a reading, not a test. Papers
proposing such fixes should report a seed-rewriting row (lineage off)
next to the usual one, compare against a no-feedback prior, and report
per-category coverage of failures the proposer was not shown. Without
the seed row, a method that works by substituting for lineage and a
method that adds to it look the same. The same holds for comparisons
between optimizers: on Banking77, one sentence in the reflection prompt
was worth more than GEPA's entire search machinery.

\textbf{Where the tokens go.} GEPA's extra cost is in validation, not in
thinking. Reflection is 2--15\% of its tokens; the rest is full
re-evaluations and minibatch checks, and an early-stop rule does not
remove the gap. Efficiency work on reflective optimizers should target
the validation schedule first.

\hypertarget{limitations}{%
\section{Limitations}\label{limitations}}

\begin{itemize}
\tightlist
\item
  \textbf{The carry-facts fix is post hoc in design.} It was formulated
  after seeing where GEPA won on Banking77, and pre-registered only
  before its own runs. It is tested on two tasks, helps on one and
  slightly hurts on the other. Whether it helps on other
  vocabulary-bound tasks is untested.
\item
  \textbf{Task construction.} Five of the main families are constructed
  so that the needed rules cannot be inferred from the task statement.
  On the two public benchmarks the lineage effect is smaller (+0.088),
  and the disclosure contrasts do not replicate (§4.8).
\item
  \textbf{Benchmark preprocessing.} Banking77 uses 12 of 77 intents (20
  gave at most 0.15 within 8 evaluations in a smoke run). BBH drops
  inputs longer than the 300-character feedback window (198 of 250
  kept). Both choices were fixed before the pre-registration, and smoke
  runs are excluded.
\item
  \textbf{Single judge.} Free-text replies use one LLM judge
  (gpt-4.1-mini, temperature 0) with per-example rubrics.
\item
  \textbf{Weak proxy for diversity.} ``Diverse'' uses Levenshtein
  farthest-point selection, a weak proxy for semantic categories; the
  clustered arm covers this gap.
\item
  \textbf{Two model families, reduced coverage.} The full grid uses
  gpt-4.1-mini. The claude-haiku-5-5 replication covers five of the
  seven tasks, with the first-k cells of its diverse arm taken from an
  earlier batch. The strong-proposer arm covers three tasks. The replies
  judge is gpt-4.1-mini throughout, and the clusterer is always the
  proposer's model.
\item
  \textbf{IFBench not run.} Its verifier requires executing external
  code, which was not approved in this study.
\item
  \textbf{ScoNe near ceiling.} Held-out is 0.89--0.92 in every arm, so
  its null result cannot separate no effect from no headroom.
\item
  \textbf{Clustered arm.} It implements only ESPO's Diagnose step, not
  Diversify or Stabilize. Its pre-registered prompt did not require the
  expected answer; the gold-keeping variant was added after seeing
  results (registered before running) and is reported as a sensitivity
  analysis.
\item
  \textbf{Batch drift.} The sensitivity, extension and fix arms ran as
  separate batches; drift in the API models between batches cannot be
  excluded.
\item
  \textbf{Lost credit.} Three batches lost API credit mid-run. Each was
  quarantined and rerun in full under a rule written before its results
  were seen (Appendix A); the quarantined runs are archived, not
  analysed.
\end{itemize}

\hypertarget{pre-registered-tests}{%
\section{Pre-registered tests}\label{pre-registered-tests}}

The pre-registrations were written in Chinese in
\texttt{09-experiments.md} before each batch of runs, timestamped, and
only appended to afterwards; each protocol was also registered with its
success threshold in \texttt{experiments.json}. Below is an English
translation of the tests. Where the translation and the original differ,
the original governs. All intervals are 95\% task-stratified bootstrap
intervals (10,000 resamples of runs within each cell), and tasks are
weighted equally.

\textbf{A.1 Main grid (2026-10-27; code frozen at git d5ab8923 plus
\texttt{results/code/}, SHA256SUMS).} Model gpt-4.1-mini for task,
proposer and clusterer (and for the replies judge, temperature 0). Seeds
0--9. Budget 8 evaluations for prompt tasks and 12 for code tasks.
Metric: held-out score of the best search-set candidate; held-out data
is never used for selection. Cells: \texttt{seed-none} and \{seed,
parent\} × \{first, diverse, minibatch, cluster\} on seven tasks.

\begin{itemize}
\tightlist
\item
  \emph{C2 (primary).} For X = diverse and X = cluster, I\_X = mean over
  tasks of (X − first)\_seed − (X − first)\_parent. Supported if the
  lower bound of I\_diverse \textgreater{} 0. R1 refutes if both
  I\_diverse and I\_cluster contain 0. R2 (collapse) if cluster −
  max(other disclosures) has lower bound \textgreater{} 0 in both rows.
\item
  \emph{C1 (crowding out).} In the seed row, coverage of unshown failure
  kinds, \texttt{seed-full-first} vs \texttt{seed-none}. Supported if
  the upper bound of (first − none) \textless{} 0 on at least two of
  labels, duration and size. R3: if every seed-row full-* minus
  seed-none contains 0 on ScoNe, C1 is narrowed to rules that can only
  be learned from feedback (ScoNe was flagged as near ceiling before
  running).
\item
  \emph{C3 (GEPA tokens).} Median tokens-to-solve, GEPA with early
  stopping vs \texttt{parent-full-first}. R4: if the ratio CI contains 1
  on at least 3 of 5 tasks, C3 is restated as ``the gap comes from
  missing early stopping''.
\end{itemize}

\textbf{A.2 Sensitivity arm \texttt{sens-clustergold} (2026-10-27,
registered after round-1 review and before running; code in
\texttt{results/code\_round1/}).} Only the cluster prompt changes: the
representative failure now carries input, wrong answer and full expected
answer. The pre-registered cluster results are kept and not replaced.

\begin{itemize}
\tightlist
\item
  \emph{S1 (information loss).} On replies, parent clustergold − cluster
  has lower bound \textgreater{} 0, and at least 80\% of the summaries
  keep the expected answer.
\item
  \emph{S2.} Parent clustergold − first on seven tasks: upper bound
  \textless{} 0 means summaries are harmful even with the target kept;
  an interval containing 0 means indistinguishable; lower bound
  \textgreater{} 0 favours full-failure diagnosis.
\end{itemize}

\textbf{A.3 Second-model replication \texttt{rep-haiku} (2026-10-27).}
claude-haiku-5-5 as proposer and task model (default adaptive thinking);
replies judge unchanged. Tasks: labels, routing, replies, duration,
size. Cells: \texttt{seed-none}, \{seed, parent\} × \{first, minibatch,
clustergold\}; n = 10.

\begin{itemize}
\tightlist
\item
  \emph{H1.} Lineage main effect (mean of the three columns), lower
  bound \textgreater{} 0.
\item
  \emph{H2.} I\_clustergold, lower bound \textgreater{} 0.
\item
  \emph{H3.} Parent clustergold − first, upper bound \textless{} +0.05.
\item
  Any H that fails is reported as ``does not hold on Haiku'', without
  added cells.
\end{itemize}

\textbf{A.4 Extension (2026-10-28T16:05 −07:00; code in
\texttt{results/code\_round4/}, data hashes in
\texttt{DATA\_SHA256SUMS}).} Design choices fixed after smoke runs and
before registration, with smoke runs excluded from analysis: Banking77
reduced from 20 to 12 intents because 20 gave at most 0.15 within 8
evaluations; for BBH salient-translation-error detection, the shared
instructions and options were moved into the seed prompt, inputs longer
than the 300-character feedback limit were dropped (198 of 250 kept),
and the search set was raised from 2 to 3 per class after the
2-per-class smoke run overfit (search 1.0, held-out 0.50). One smoke run
showed Banking77 parent clustergold 0.22 vs parent first 0.72 (n = 1);
the tests were not changed.

\begin{itemize}
\tightlist
\item
  \emph{\texttt{public-bench}.} gpt-4.1-mini on banking77 and bbh\_ste,
  nine cells (\texttt{seed-none} and \{seed, parent\} × \{first,
  diverse, minibatch, clustergold\}) run interleaved. P1 lineage
  (four-column mean), lower bound \textgreater{} 0. P2 I\_clustergold,
  lower bound \textgreater{} 0. P3 parent clustergold − first, upper
  bound \textless{} +0.05. P4 I\_diverse, lower bound \textgreater{} 0.
  P5 parent minibatch − first, upper bound \textless{} 0. A task with
  seed-none ≥ 0.90 is treated as at ceiling and every test is reported
  with and without it. A failed test is reported as ``does not hold on
  public benchmarks''.
\item
  \emph{\texttt{gepa-public}.} GEPA 0.0.7 (minibatch 3, with and without
  early stopping), 10 runs per task, against \texttt{parent-full-first}.
  R4 holds on a task if both arms solve at least 3/10 and the ratio CI
  contains 1.
\item
  \emph{\texttt{strong-proposer}.} gpt-4.1 as proposer and clusterer,
  gpt-4.1-mini as task model and judge; labels, routing, replies; cells
  \texttt{seed-none}, \{seed, parent\} × \{first, clustergold\}. S1
  lineage (mean of the two columns), lower bound \textgreater{} 0; S2
  I\_clustergold, lower bound \textgreater{} 0; S3 parent clustergold −
  first, upper bound \textless{} +0.05.
\item
  \emph{\texttt{rep-haiku-diverse}.} claude-haiku-5-5, \{seed,
  parent\}-full-diverse on the five replication tasks. Hd: I\_diverse,
  lower bound \textgreater{} 0, using the first-k cells from
  \texttt{rep-haiku} (not interleaved).
\end{itemize}

\textbf{A.5 Fix arm \texttt{fix-carryfacts} (2026-10-28, written in
\texttt{09-experiments.md} after the diagnosis of §4.3 and before any
fix run; code in \texttt{results/code\_round5/}).} One sentence
(\texttt{CARRY\_FACTS}, quoted in §4.3) is appended to the proposer's
objective; nothing else changes. gpt-4.1-mini in all roles; banking77
and bbh\_ste; cells \texttt{seed-full-first} and
\texttt{parent-full-first}; n = 10; compared against
\texttt{parent-full-first} and GEPA mb3 from A.4 (not interleaved).
Because the fix was designed after seeing where GEPA won, the paper
labels it post hoc in design.

\begin{itemize}
\tightlist
\item
  \emph{F1.} Banking77 parent-carry − parent-first, lower bound
  \textgreater{} 0.
\item
  \emph{F2.} Banking77 GEPA − parent-carry: an upper bound ≤ 0 means the
  gap is explained by the meta-prompt; otherwise not fully explained.
\item
  \emph{F3 (no harm).} BBH parent-carry − parent-first, lower bound
  \textgreater{} −0.05.
\item
  \emph{F4.} Lineage under carry-facts (parent-carry − seed-carry, mean
  of both tasks), lower bound \textgreater{} 0.
\item
  Outcome: F1, F2 and F4 hold; F3 fails (−0.061 {[}−0.119, −0.006{]})
  and is reported in §4.8.
\end{itemize}

  \bibliography{../refs.bib}

\end{document}
